\documentclass[12pt]{article}
\RequirePackage[brazilian]{babel}
\RequirePackage{amsmath}
\RequirePackage{hyperref}
\RequirePackage[top=0cm, bottom=2cm, left=2cm, right=2cm]{geometry}

\title{\textbf{Simulador de Circuitos Elétricos em C}}
\author{} % Evita o aviso de autor ausente
\date{}

\begin{document}

\maketitle

\section*{Resumo \& Visão Geral}
Este projeto une Estrutura de Dados (Grafos) com Teoria de Circuitos Elétricos (Análise Nodal) para simular o comportamento de circuitos de corrente contínua (CC) em linguagem C.

\section{1. Mapeamento: Visão Elétrica \texorpdfstring{$\rightarrow$}{->} Netlist \texorpdfstring{$\rightarrow$}{->} Grafo}

A conversão do circuito físico para a memória do programa segue este fluxo visual:

\subsection*{A. Circuito Elétrico (Desenho)}
\begin{verbatim}
       (Nó 1) ------------ (Nó 2)
         |                   |
       [10V]               [100 Ohm]
        (V1)                (R1)
         |                   |
       (Nó 0) ------------ (Nó 0)  <-- Ref (GND)
\end{verbatim}

\subsection*{B. Arquivo Netlist (\texttt{circuito.txt})}
\begin{verbatim}
V1 1 0 10.0
R1 1 2 100.0
R2 2 0 50.0
\end{verbatim}

\subsection*{C. Estrutura de Grafo em Memória (C)}
\begin{verbatim}
  [ Nó 0 (GND) ] <==== V1 (10V) ====> [ Nó 1 ]
        ||                               ||
     R2 (50 Ohm)                     R1 (100 Ohm)
        ||                               ||
        +================================+
\end{verbatim}

---

\section{2. Arquitetura das Estruturas de Dados}

\begin{itemize}
    \item \textbf{Vértice (Nó):} Guarda a tensão calculada ($V_{\text{nó}}$) e a lista de componentes conectados.
    \item \textbf{Aresta (Componente):} Guarda o tipo (\texttt{R}, \texttt{V}, \texttt{I}), o valor e os nós \texttt{A} e \texttt{B}.
\end{itemize}

---

\section{3. Motor de Simulação (\texorpdfstring{$A \cdot X = B$}{A . X = B})}

Transformação das leis de Kirchhoff na matriz nodal resolvida por Eliminação Gaussiana:

\begin{verbatim}
  Matriz Condutâncias (A)     Vet. Tensões (X)    Vet. Fontes (B)
   [  G11   -G12   ...  ]  *   [   V_1   ]     =    [   I_1   ]
   [ -G21    G22   ...  ]      [   V_2   ]          [   I_2   ]
\end{verbatim}

---

\section{4. Roteiro Prático de Desenvolvimento}

\begin{enumerate}
    \item \textbf{Leitura do Arquivo:} Ler linha a linha do \texttt{.txt} e instanciar os componentes.
    \item \textbf{Montagem do Grafo:} Conectar as pontas dos componentes aos nós correspondentes.
    \item \textbf{Estampagem:} Preencher a matriz $A$ com as condutâncias ($G = 1/R$).
    \item \textbf{Cálculo:} Resolver o sistema de equações e exibir o relatório final no terminal.
\end{enumerate}

\end{document}